\section{Related Work}
\label{sec:related}

\paragraph{Video-language models.}
The recent wave of video-language models pairs a visual encoder with a large language model.
Video-ChatGPT~\cite{maaz2023videochatgpt} averages spatial and temporal features and fine-tunes on video-instruction data.
VideoLLaVA~\cite{lin2023videollava} adds a projection layer that aligns video and image representations before the language model.
LLaVA-Next-Video~\cite{zhang2024llavanextvideo} scales frame resolution and count.
SeViLA~\cite{yu2023sevila} localizes keyframes first, then answers questions from those frames.
Chat-UniVi~\cite{jin2024chatunivi} merges spatial and temporal features at multiple granularities with adaptive tokens.
BLIP~\cite{li2022blip} works at the single-image level, training one model for captioning and VQA. We adopt BLIP as the perception backbone in \ours.
Each of these systems expects the full video at inference time. A growing stream with no known end breaks that expectation. We built \ours for exactly that regime: causal, frame-by-frame input with bounded memory.

\paragraph{Streaming perception.}
Yang \etal~\cite{yang2022real} formalize real-time streaming perception for object detection. Their benchmark penalizes a method whenever it falls behind the input frame rate.
StreamPETR~\cite{wang2023streampetr} carries this idea to 3D detection by propagating temporal features across frames.
Several recent efforts come closer to our problem.
VideoStreaming~\cite{qian2024streaming} proposes a training protocol aware of streaming constraints.
Flash-VStream~\cite{zhang2024flashvstream} adds a memory layer for long video streams.
VideoLLM-online~\cite{chen2024videollmonline} extends large language models with a streaming decoding mechanism so the model can answer while the video is still playing.
Dispider~\cite{he2024dispider} disentangles perception, decision, and reaction into parallel modules, enabling active real-time interaction with video feeds.
Despite this progress, none of these methods explicitly reason about which temporal window a user's question refers to. \ours fills that gap with temporal scope classification and scope-aware memory filtering, achieving stronger results on both established and new streaming benchmarks.

\paragraph{Video question answering.}
NExT-QA~\cite{xiao2021nextqa} tests causal and temporal reasoning on short clips and reveals an important gap: models that excel at descriptive questions struggle on causal and temporal ones.
Its per-type evaluation (causal, temporal, descriptive) inspired our per-scope evaluation (instant, recent, historical), and its open-ended QA results show that strong multi-choice models struggle to \emph{generate} correct answers---a challenge \ours inherits through its open-ended Flan-T5 generation stage.
STAR~\cite{wu2021star} probes situated reasoning with compositional questions, and MVBench~\cite{li2024mvbench} covers a broad range of video understanding tasks.
EgoSchema~\cite{mangalam2023egoschema} pushes the time horizon to minutes-long egocentric clips from Ego4D~\cite{grauman2022ego4d}, a large-scale egocentric video dataset spanning thousands of hours of daily activities.
MovieChat~\cite{song2024moviechat} tackles long videos by merging frame features into a fixed-size buffer.
OVO-Bench~\cite{li2025ovobench} is the closest antecedent: it evaluates Video-LLMs on online video understanding with questions that must be answered causally, without seeing future frames.
We evaluate on OVO-Bench alongside NExT-QA and EgoSchema to cover a wide spectrum.
LiveQA-Bench complements these benchmarks by introducing explicit temporal scope labels, enabling per-scope analysis that reveals \emph{which} components matter for \emph{which} temporal reasoning demands.

\paragraph{Two-stage pipelines and temporal memory.}
Separating perception from language generation lets each stage use a frozen model, removing the need for joint training.
BLIP-2~\cite{li2023blip2} pairs a frozen image encoder with Flan-T5~\cite{chung2022flanT5} to answer visual questions without task-specific training; our SFG follows the same two-stage idea at the frame level.
For temporal memory, Recurrent Memory Transformers~\cite{bulatov2022recurrent} persist memory tokens across segments, while Memorizing Transformers~\cite{wu2022memorizing} use a $k$-NN index over an external bank.
In the video domain, MovieChat~\cite{song2024moviechat} merges frame features into a buffer and Ma \etal~\cite{ma2023livelylivechat} compress tokens for long conversations.
\ours departs from these designs by ranking stored frames with a visual-semantic importance score (not recency alone) and adding a temporal router that selects only the time slice the question refers to.

\begin{figure*}[!t]
  \centering
  \begin{tikzpicture}[
    node distance=0.6cm and 1.2cm,
    box/.style={draw, rounded corners=4pt, minimum height=0.9cm, minimum width=2.0cm, align=center, font=\small},
    io/.style={draw, rounded corners=4pt, minimum height=0.8cm, minimum width=1.6cm, align=center, font=\small\itshape, dashed},
    arr/.style={-{Stealth[length=5pt]}, thick},
    lbl/.style={font=\scriptsize, midway, fill=white, inner sep=1pt, yshift=3pt},
  ]
    % Row 1: video stream path
    \node[io] (stream) {Live Video\\Stream};
    \node[box, right=1.2cm of stream, fill=blue!8] (enc) {CLIP\\Encoder};
    \node[box, right=1.0cm of enc, fill=blue!8] (skm) {Semantic\\Keyframe\\Memory};
    \node[box, right=1.6cm of skm, fill=green!8] (blip) {BLIP\\Perception};
    \node[box, right=1.4cm of blip, fill=green!8] (flan) {Flan-T5\\Synthesis};
    \node[io, right=1.2cm of flan] (answer) {Real-Time\\Answer};

    % Row 2: query path
    \node[io, below=1.6cm of stream] (query) {User\\Question};
    \node[box, right=1.2cm of query, fill=orange!10] (tqr) {Temporal Query\\Router};

    % Arrows: video stream path
    \draw[arr] (stream) -- (enc) node[lbl, above] {frames};
    \draw[arr] (enc) -- (skm) node[lbl, above] {$v_t$};
    \draw[arr] (skm) -- (blip) node[lbl, above] {keyframes};
    \draw[arr] (blip) -- (flan) node[lbl, above] {captions};
    \draw[arr] (flan) -- (answer) node[lbl, above] {$a$};

    % Arrow: query to TQR
    \draw[arr] (query) -- (tqr) node[lbl, above] {$q$};

    % Arrow: TQR to BLIP via clean L-path (right then up)
    \coordinate (turn) at (blip.south |- tqr.east);
    \draw[arr] (tqr.east) -- (turn) node[lbl, above, pos=0.7] {scope} -- (blip.south);

    % Dashed module groups
    \begin{scope}[on background layer]
      \node[fit=(enc)(skm), draw, dashed, rounded corners=6pt, inner sep=8pt,
            label={[font=\scriptsize]above:SKM}] {};
      \node[fit=(tqr), draw, dashed, rounded corners=6pt, inner sep=8pt,
            label={[font=\scriptsize]above:TQR}] {};
      \node[fit=(blip)(flan), draw, dashed, rounded corners=6pt, inner sep=8pt,
            label={[font=\scriptsize]above:SFG (two-stage)}] {};
    \end{scope}
  \end{tikzpicture}
  \caption{\textbf{Architecture of \ours.}}
  \label{fig:architecture}
\end{figure*}

